\documentclass[12pt,a4paper]{article}
\usepackage[UTF8]{ctex}
\usepackage{geometry}
\usepackage{amsmath,amssymb}
\usepackage{enumitem}
\usepackage{fancyhdr}
\usepackage{setspace}
\usepackage{hyperref}

\geometry{left=2.5cm,right=2.5cm,top=2.5cm,bottom=2.5cm}
\setlength{\headheight}{15pt}
\setstretch{1.5}
\raggedbottom
\setlength{\emergencystretch}{3em}
\hfuzz=100pt
\hbadness=10000
\hypersetup{hidelinks}

\pagestyle{fancy}
\fancyhf{}
\fancyhead[C]{说明书}
\fancyfoot[C]{\thepage}

\begin{document}

\begin{center}
    {\Huge\heiti 说明书}
\end{center}

\vspace{1em}

\noindent\textbf{发明名称：}量子-经典异构光纤网络的动态波分复用与路由控制方法

\section*{一、技术领域}

本发明涉及量子保密通信、光传送网、波分复用、软件定义网络控制以及量子网络资源编排技术领域，尤其涉及一种面向量子信号与经典高速业务信号共纤传输场景的动态波分复用与路由控制方法、系统、设备及存储介质。

更具体地，本发明适用于量子密钥分发网络与现网 DWDM、OTN、ROADM 系统协同部署的场景，用于在同一根光纤中同时承载量子信号与经典业务信号时，对量子信道的路由、波长、发射时间窗、探测门控参数和纠错配置进行联动控制。

\section*{二、背景技术}

未来量子保密通信网络难以依赖大规模专属暗光纤进行部署。在运营商现网中，更现实的路径是在既有光传送网上实现量子信号与经典业务信号的共纤传输。然而，量子信号通常处于单光子或弱相干光功率等级，而经典业务信号功率远高于量子信号，二者在同纤传输时会引入显著的拉曼散射噪声、串扰、滤波泄漏以及链路预算波动，进而导致量子误码率升高、安全密钥率下降甚至链路失效。

现有技术中，一类方案主要围绕量子密钥分发链路与经典数据链路的物理层共存展开，通过调节信道间隔、增加滤波器、优化功率谱等方式减小噪声影响；另一类方案将量子业务融入光网络资源分配框架，采用静态或离线的路由与波长分配方法实现量子信道承载；还有一类方案采用时间调度策略，在特定时隙内发送量子信号以降低共纤噪声背景。

但是，现有方案普遍存在以下不足：
\begin{enumerate}[leftmargin=2em]
    \item 大多面向静态规划、离线资源分配或实验室级链路验证，缺少针对运行期波动的在线闭环控制；
    \item 往往仅关注路由或波长中的单一维度，缺少对路由、波长、时间窗、探测门控与纠错配置的联合优化；
    \item 缺少面向 ROADM、OTN 和 QKD 设备的统一执行编排，难以直接落地到现网控制平面；
    \item 对短时尖峰噪声、遥测异常、设备重构失败和控制震荡等工程问题考虑不足；
    \item 在执行波长迁移或路径迁移时，缺少新旧配置重叠保护和回滚机制，可能造成量子业务中断。
\end{enumerate}

因此，亟需一种能够感知运行期经典噪声环境并进行在线联合决策的动态波分复用与路由控制方案，以提高量子业务在现网共纤部署中的稳定性、可用性和工程可实施性。

\section*{三、发明内容}

\subsection*{（一）发明目的}

本发明的目的在于提供一种量子-经典异构光纤网络的动态波分复用与路由控制方法、系统、设备及存储介质，以解决现有技术中量子与经典业务共纤场景下缺少在线联合优化、运行期闭环控制以及平滑重构保护机制的问题。

\subsection*{（二）技术方案}

为实现上述目的，本发明提供如下技术方案。

一种面向量子-经典异构光纤网络的动态波分复用与路由控制方法，应用于量子控制器或编排器，所述异构光纤网络在同一光纤中承载量子信号与经典波分复用业务信号，并与 ROADM、OTN 设备及 QKD 设备协同；所述方法包括：
\begin{enumerate}[leftmargin=2em]
    \item 采集链路段级的经典功率谱、波长占用信息、放大器状态和告警信息，以及量子链路测量指标；
    \item 基于所述采集数据构建链路段、波长和时间窗三级噪声感知数据，形成三维噪声张量，并计算候选路径、候选量子波长以及候选发射时间窗在给定探测门控参数和纠错配置下的预测量子误码率与预测安全密钥率；
    \item 在满足量子业务质量约束和现网扰动约束的条件下，对路由、量子波长、发射时间窗、纠错配置和探测门控参数进行联合优化，得到目标决策；
    \item 按先轻后重的分层执行策略下发重配置，优先执行发射时间窗、门控参数和纠错配置的轻量调整，在轻量调整后仍不满足质量约束时，再执行量子波长迁移或路径迁移；
    \item 在执行迁移前生成上一稳定配置快照，并在迁移窗口内保持新旧配置重叠保护；当重配置失败或者重配置后量子业务质量未达标时，执行回滚并记录审计信息。
\end{enumerate}

在一个实施方式中，所述量子链路测量指标至少包括量子误码率 $QBER$、单光子计数率、暗计数率、探测器饱和度以及安全密钥率 $R_{\mathrm{sec}}$。所述三维噪声张量可表示为
\[
\mathcal{N}(p,\lambda,t),
\]
其中 $p$ 表示候选路径，$\lambda$ 表示候选量子波长，$t$ 表示候选发射时间窗。

在一个实施方式中，系统基于所述三维噪声张量、占用张量以及器件状态，对候选决策进行预测评估。可通过如下形式的目标函数进行联合优化：
\[
(p^\star,\lambda^\star,t^\star,e^\star,g^\star)=\arg\max U,
\]
其中
\[
U=\alpha R_{\mathrm{sec}}-\beta QBER-\gamma C_{\mathrm{reconf}}-\delta C_{\mathrm{impact}},
\]
$e$ 表示纠错配置，$g$ 表示探测门控参数，$C_{\mathrm{reconf}}$ 表示重构代价，$C_{\mathrm{impact}}$ 表示对现网经典业务的扰动代价。

在一个实施方式中，所述分层执行策略至少包括：
\begin{enumerate}[leftmargin=2em]
    \item L1 协议层轻量动作：调整发射时间窗、探测门控宽度、门控相位、诱骗态比例以及纠错码率和块长；
    \item L2 光层中量动作：在不改变端到端拓扑的前提下执行量子波长迁移，并与滤波或保护带策略同步切换；
    \item L3 路由级重构动作：执行路径迁移，必要时联合量子波长迁移。
\end{enumerate}

在一个实施方式中，系统设置切换滞回和最小驻留时间，仅当预测劣化趋势持续预设采样周期数且效用增益超过阈值时，才允许从较轻层级升级到较重层级，以避免频繁切换造成控制震荡。

在一个实施方式中，当遥测缺失、数据异常或链路状态可信度不足时，系统进入保守模式，冻结光层重构，仅允许执行轻量动作，并提高触发门限。

在一个实施方式中，当判定当前噪声劣化属于短时尖峰时，不执行波长迁移和路径迁移，而通过时间窗错避、门控收窄和纠错增强完成快速抑制。

本发明还提供一种面向量子-经典异构光纤网络的动态波分复用与路由控制系统，包括：
\begin{enumerate}[leftmargin=2em]
    \item 经典噪声遥测模块，用于采集 ROADM、OTN 的功率谱、波长占用、端口功率、放大器状态和告警信息；
    \item 量子链路测量模块，用于采集 QKD 收发端的量子误码率、计数率、暗计数率、探测门控状态、饱和度和密钥率；
    \item 噪声建模模块，用于构建链路段、波长和时间窗三级噪声张量并输出预测指标；
    \item 联合优化引擎，用于对路由、波长、时间窗、纠错配置和门控参数进行联合求解并输出重构计划；
    \item 重配置编排器，用于向 ROADM、WSS、OTN 设备和 QKD 设备下发重配置指令，并执行重叠保护与回滚；
    \item 策略审计模块，用于记录调度依据、阈值、影响评估和执行结果。
\end{enumerate}

本发明还提供一种电子设备以及一种计算机可读存储介质，用于执行上述方法。

\subsection*{（三）有益效果}

与现有技术相比，本发明至少具有如下有益效果：
\begin{enumerate}[leftmargin=2em]
    \item 实现链路段、波长和时间窗三级噪声感知，能够更准确地表征运行期共纤噪声环境；
    \item 实现路由、波长、时间窗、门控和纠错配置的联合优化，而非单维资源调整；
    \item 支持“先协议层、后光层、再路由”的分层控制，在兼顾现网稳定性的同时降低量子业务失效风险；
    \item 通过新旧配置重叠保护、可回滚快照和短期禁用机制提高工程可落地性；
    \item 通过统一的南向和北向接口实现与 ROADM、OTN、QKD 设备及密钥管理系统的闭环协同；
    \item 在不新增暗光纤的前提下，提高量子业务在动态经典业务负荷下的持续可服务性。
\end{enumerate}

\section*{四、附图说明}

为了更清楚地说明本发明实施例或现有技术中的技术方案，下面对实施例中所需要使用的附图作简单说明。所述附图仅为示意性附图。

\begin{enumerate}[leftmargin=2em]
    \item 图1为本发明量子控制器与 ROADM、OTN、QKD 设备协同的系统架构示意图；
    \item 图2为本发明链路段、波长和时间窗三级噪声感知与联合优化流程示意图；
    \item 图3为本发明 L1、L2 和 L3 分层执行及回滚机制示意图；
    \item 图4为本发明新旧配置重叠保护与执行回滚时序示意图；
    \item 图5为本发明在城域环网场景中进行量子波长迁移和路径迁移的实施例示意图。
\end{enumerate}

\section*{五、具体实施方式}

下面结合附图和具体实施方式对本发明作进一步详细说明。应当理解，这里描述的实施方式仅用于解释本发明，并不用于限定本发明。

\subsection*{（一）术语定义}

在本说明书中：
\begin{itemize}[leftmargin=2em]
    \item ``候选路径''记为 $p$，表示可供量子业务选择的端到端或分段承载路径；
    \item ``量子波长''记为 $\lambda$，表示用于量子信号传输的候选波长通道；
    \item ``发射时间窗''记为 $t$，表示用于发送量子信号的离散时隙；
    \item ``纠错配置''记为 $e$，表示纠错码率、块长、交织参数等纠错相关参数；
    \item ``探测门控参数''记为 $g$，表示探测门控宽度、门控相位、门控偏置等接收相关参数；
    \item ``三维噪声张量''记为 $\mathcal{N}(p,\lambda,t)$，表示候选路径、波长和时间窗下的噪声环境描述。
\end{itemize}

\subsection*{（二）系统架构实施方式}

如图1所示，本发明实施例中的系统至少包括量子控制器 100、ROADM/WSS 节点 200、OTN 交叉节点 300、QKD 收发设备 400 以及密钥管理或运维系统 500。

量子控制器 100 用于接收来自 ROADM/WSS 节点 200 和 OTN 交叉节点 300 的功率谱、占用和告警信息，同时接收来自 QKD 收发设备 400 的量子链路测量指标。量子控制器 100 内部可进一步包括经典噪声遥测模块、量子链路测量模块、噪声建模模块、联合优化引擎、重配置编排器和策略审计模块。

在一个实施方式中，经典噪声遥测模块通过 NETCONF、gNMI、SNMP 或网管北向接口获取功率谱和端口状态信息；量子链路测量模块通过量子设备控制或遥测接口获取量子误码率和密钥率信息；重配置编排器通过 SDN 控制器或网管接口下发对 ROADM、WSS 和 OTN 的重构指令。

\subsection*{（三）三级噪声感知与联合优化实施方式}

如图2所示，在一个实施方式中，系统以预设采样周期滚动收集链路段级经典业务占用信息和量子链路测量信息，并形成占用张量与噪声张量。对任一候选组合 $(p,\lambda,t,e,g)$，系统估计其预测量子误码率和预测安全密钥率。

在一个实施方式中，若某一链路段在特定时间窗内新开通高功率经典业务，则该时间窗对应的噪声张量分量将上升，系统据此可优先尝试执行时间窗错避或门控收窄，而不是直接触发路径迁移。

在一个实施方式中，联合优化引擎可采用启发式算法、整数规划、两阶段候选筛选加在线快速选择、强化学习或模型预测控制实现。无论采用何种具体求解器，其输入均包括候选路径集合、候选波长集合、候选时间窗集合、量子业务质量门限以及对经典业务扰动代价的约束。

\subsection*{（四）分层执行与回滚实施方式}

如图3和图4所示，本发明优选采用分层执行策略。当某一候选决策违反量子业务质量约束但尚未达到触发重构的程度时，系统优先执行 L1 轻量动作，例如调整发射时间窗、门控宽度和纠错配置；当 L1 无法使指标恢复至目标范围时，再升级到 L2 或 L3。

在一个实施方式中，执行光层动作或路由动作前，重配置编排器保存上一稳定配置快照，并启动保护窗口。在保护窗口内，新旧配置并行保持可用，只有在收到 ROADM/WSS/OTN 的执行回执且量子链路指标达到门限后，才释放旧配置。若执行失败或执行后的量子误码率未恢复，则在保护窗口内回滚至上一稳定配置。

在一个实施方式中，为避免反复尝试同一失败动作，系统将失败动作标记为短期禁用，并在预设时长内不再选择。

\subsection*{（五）异常工况处理实施方式}

在一个实施方式中，当检测到遥测缺失、监测结果冲突或测量可信度不足时，系统进入保守模式，冻结 L2 和 L3 动作，仅保留 L1 轻量调整，并触发补采样。

在一个实施方式中，当检测到短时噪声尖峰时，系统优先执行时间窗错避、门控收窄和纠错增强，而不触发波长迁移或路径迁移。

在一个实施方式中，当探测器计数率或饱和度达到阈值时，系统优先降低发射占空比、缩短门控宽度或调整诱骗态配置，并在必要时切换至更低噪声波长。

\subsection*{（六）城域环网实施例}

在一个城域环网实施例中，如图5所示，8 个 ROADM 节点构成环网，网络中承载多个经典 DWDM 信道和量子信道。控制器检测到晚高峰期间某一路段 400G 业务增开，导致当前量子波长所在频段的预测拉曼噪声上升。系统首先执行 L1 动作，缩窄探测门控并调整发射时间窗；当轻量调整后预测安全密钥率仍低于阈值时，系统触发 L2 动作，将量子业务迁移至备用量子波长；若迁移后仍不满足要求，则进一步执行 L3 动作，绕开噪声链路段完成路径迁移。整个过程中，新旧配置在保护窗口内并行保持可用，若任一步骤失败则回滚至上一稳定配置。

\subsection*{（七）适用范围说明}

本发明不限于 BB84 协议，也可适用于诱骗态 BB84、DPS-QKD、CV-QKD 及其他量子密钥分发协议；不限于 DWDM、OTN 或 ROADM 场景，也可扩展至 PON 或其他具有量子与经典共纤需求的承载网络。

本发明亦不限于单一控制器部署方式，可由集中式控制器、边缘控制器或分层控制平面实现。只要具备所述三级噪声感知、联合优化、分层执行与回滚保护的基本技术思想，均应落入本发明的保护范围。

\end{document}
